\begin{lemma}\label{lem:nonic242-quotient5}
Put $X=X_1(2,42)$, and let $G_9$ be generated by
\[
 a(P,Q)=(21Q,P+Q),\qquad b(P,Q)=(P,25Q).
\]
Then $G_9\simeq C_3\times C_3$ is defined over $\mathbb Q$.
The quotient $Y=X/G_9$ has genus $5$ and the following function
field. Put
\[
 t=u^2,\quad D=t^2+13t+49,\quad N=t^2+5t+1,\quad
 M=t^4+14t^3+63t^2+70t-7.
\]
Adjoin the coordinates $(x_3,y_3)$ of a nonzero point of order $3$
on
\[
 E_u:\quad y^2=x^3-\frac{N}{48D}x+\frac{M}{864D^2}.
\]
Thus $x_3$ satisfies the $3$-division polynomial, and
$y_3^2=x_3^3-Nx_3/(48D)+M/(864D^2)$.
The Jacobian $J_Y(\mathbb Q)$ has rank zero.
\end{lemma}
\begin{proof}
The two displayed transformations commute and have order $3$.
The first cyclically permutes the three nonzero points of order $2$.
The second fixes the full $2$-structure and multiplies the point
of order $7$ by $4$. They generate a group of order $9$.

We give the descent calculation. The Tate normal form for a point
of order $7$ has
\[
 b_7=v^2(v-1),\qquad c_7=v(v-1).
\]
Put $H=v^2-v+1$, $s=(v^3-3v+1)/(v(v-1))$, and $t=s-8$.
The diamond operator of order $3$ sends $v$ to $(v-1)/v$.
Completing the square and depressing the cubic, followed by
scaling by $H$, gives $E_u$ after adjoining $u^2=t$.
The exact invariant identities are
\[
 D=\frac{H^3}{v^2(v-1)^2},\quad
 \frac{c_4}{H^4}=\frac ND,\quad
 \frac{c_6}{H^6}=-\frac M{D^2},\quad
 \frac{\Delta}{H^{12}}=\frac{t}{D^4}.
\]
The square root $u$ records the orientation of the full
$2$-structure, and is fixed by its alternating subgroup.
The point $14Q$ supplies the nonzero point of order $3$.
These functions are invariant under $G_9$.

The resulting field has degree $8$ over $\mathbb Q(u)$.
The two successive equations have degrees $4$ and $2$ and remain
irreducible over $\mathbb F_5(u)$.
Moreover
\[
 j=\frac{N^3D}{t}
\]
has degree $16$ as a function of $u$. Hence the field has degree
$128$ over $\mathbb Q(j)$. The projective index for $X$ is $1152$,
so this is the full invariant field of $G_9$.
The quotient has good reduction at $5$, since $X$ has good
reduction there and $\#G_9=9$ is prime to $5$.
Its normalized finite function field has genus $5$.
Finally $J_Y$ is a quotient of $J_X$, whose rational rank is zero
by \cite[Theorem~3.1]{DEHMZ}.
\end{proof}

\begin{lemma}\label{lem:nonic242-quotient5-hecke}
Normalize $J_Y$ at either rational cusp. The reduction image
of $J_Y(\mathbb Q)$ at $5$ is contained in a subgroup
\[
 K\simeq C_2\times C_2\times C_2\times C_{24}
 \subset J_Y(\mathbb F_5),
\]
of index $18$.
\end{lemma}
\begin{proof}
The curve $Y$ has eight geometric cusps. Their fields are
$\mathbb Q,\mathbb Q,\mathbb Q(\sqrt{21}),
\mathbb Q(\sqrt{-7}),\mathbb Q(\sqrt{-3})$,
where each quadratic field contributes two geometric cusps.
To distinguish the branches, put
\[
 \rho=\frac{M}{12ND},\qquad h=\frac{y_3}{x_3-\rho}.
\]
At $u=0$, the two branch pairs have residues
$h^2=1/84$ and $h^2=-1/28$.
At $u=\infty$, they have residues
$h^2=1/4$ and $h^2=-1/12$.
The first three quadratic fields and the two rational cusps follow.
At $5$ the cusp signature is $1^4 2^2$.

The proven finite Jacobian is
\[
 J_Y(\mathbb F_5)=[2,6,6,48].
\]
The finite cusps generate its entire $2$-primary subgroup,
with invariants $[2,2,2,16]$.
The global rational cuspidal divisors have image $[2,2,8]$.
We do not assume that this is the entire rational image.

On $Y$, the diamond operator $\langle11\rangle$ sends
$(u,x_3,y_3)$ to $(u,x_3,-y_3)$.
The operator
\[
 A_{11}=T_{11}-1-11\langle11\rangle
\]
annihilates $J_Y(\mathbb Q)$.
Indeed the Jacobian has rank zero. Its rational torsion injects
at $5$, and has no $11$-primary part since
$\#J_Y(\mathbb F_5)=3456$. Reduction at $11$ is therefore
injective on this torsion, and the Eichler--Shimura relation gives
the assertion.

We compute $A_{11}$ on the finite cusps using their exact global
Frobenius and diamond permutations. Three ordinary closed cubic
places, together with the finite cusps, generate
$J_Y(\mathbb F_5)$. For each cubic place we lift the level
structure over $\mathbb F_{5^9}$ and compute the twelve cyclic
$11$-isogenies. Equal projected closed places are combined before
dividing the norm divisor by $3$; every resulting multiplicity
is checked to be integral.
The images respect every relation of the full finite group.
The exact kernel is $[2,2,2,24]$, and the image is $[3,6]$.
Independent Sage calculations reconstruct all $36$ checked
kernel-polynomial isogenies and verify that their norms exhaust
the respective $11$-division polynomials. They also reconstruct
the full operator and its kernel by integer linear algebra.
Thus the rational image is contained in $K=\ker A_{11}$.
\end{proof}

\begin{proposition}\label{prop:nonic242-quotient5-partial}
The two complete maps to $X_1(2,14)$ leave $919$ noncuspidal
support types for a degree-$9$ point on $X_1(2,42)$ at $5$.
The necessary condition from Lemma~\ref{lem:nonic242-quotient5-hecke}
leaves $119$ of these types. Combining it with
Lemma~\ref{lem:nonic242-even-stabilizers} leaves $63$ types.
These are necessary conditions for a degree-$9$ point.
\end{proposition}
\begin{proof}
Let $\pi:X\to Y$ be the quotient map, and let $c_0$ be the chosen
rational cusp of $Y$. For the hypothetical trace divisor $D$,
\[
 [\pi_*D-9c_0]\in J_Y(\mathbb Q).
\]
We reduce this condition at $5$ and use the containing subgroup $K$.
For every noncuspidal support, the residue-degree norm
multiplicity is included in its pushforward.
Every finite cuspidal complement of the required degree is
allowed. This is a containing set, so no global lifting assertion
for these complements is needed.
All source cases with $j=0$ are retained when the chosen coordinate
chart is not used. In particular no pole of that chart gives an
exclusion.

We test all $919$ types and obtain $119$ survivors. There are
$735$ distinct source images, of which eight have $j=0$ and are
retained. Independent Sage arithmetic checks the quotient
coordinates and residue degrees, all cusp complement classes,
all normalized class sums, and every subgroup test.
Intersecting the resulting list with the even-stabilizer
obstruction gives $63$ survivors, recorded separately.
\end{proof}

\begin{lemma}\label{lem:nonic242-primitive-cubic}
Let $D$ be a noncuspidal rational closed divisor of degree $9$ on
$X_1(2,42)$. Suppose that its reduction at $5$ is an irreducible
place $P_9$ fixed by a nonidentity $h\in G_9$ of order $3$.
Then $(\pi_*D)/3$ is a rational effective divisor of degree $3$
on $Y$. In particular, if $\pi_*P_9=3P_3$, then
\[
 [P_3-3c_0]\in K.
\]
\end{lemma}
\begin{proof}
The torsion class $[hD-D]$ reduces to zero at $5$, so it is zero.
Geometric gonality is at least $12$, and hence $hD=D$.
As in Lemma~\ref{lem:nonic242-even-stabilizers}, the automorphism
$h$ acts freely on the geometric support of $D$.
Thus $D$ descends to a closed divisor of degree $3$ on
$X/\langle h\rangle$. Pushing it to $Y$ gives $(\pi_*D)/3$.
Its normalized class is rational, and reduction gives the stated
condition by Lemma~\ref{lem:nonic242-quotient5-hecke}.
\end{proof}

The four maps to $X_1(21)$ retain $2Q$, first on $E$ and then on
each quotient of $E$ by a nonzero rational point of order $2$.
Here $J_1(21)(\mathbb Q)=[364]$ by
\cite[Theorem~4.13 and Table~2]{DEHMZ}, whereas its proven finite
Jacobian at $5$ is $[2184]$.
Testing every finite cusp complement leaves $44$ of the $119$
supports from Proposition~\ref{prop:nonic242-quotient5-partial}.
The independent replay checks $297$ normed maps, including
$175$ explicit V\'elu quotients. The optional even-stabilizer
obstruction leaves $18$ supports.

Exactly three of these supports are irreducible degree-$9$
places. All three have a diagonal order-$3$ stabilizer in $G_9$.
Lemma~\ref{lem:nonic242-primitive-cubic} removes the supports
numbered $2$ and $5$ in the complete degree-$9$ source orbit list;
the support numbered $8$ remains. Its unweighted cubic quotient
class has order $24$ and belongs to $K$.
Thus these tests leave $16$ support types, including one
irreducible degree-$9$ place. It remains to exclude these supports.
